%% Rechnungsvorlage "raster" – wird vom InvoiceRenderer gefüllt, mit denselben
%% Werten und Makroaufrufen wie "balken" (ADR-0011). Platzhalter, Makros und die
%% Regel, wer was entscheidet, beschreibt der Kopf von balken.tex. Kurz: Jede
%% einzelne Angabe darf leer sein, und diese Vorlage fragt dafür nur \ifblank.
%%
%% Gestaltungsprinzip: klassisch-tabellarisch und schwarz-weiß. Die Ordnung
%% trägt das Gitter der Positionen, nicht eine Farbe. Oben rechts der
%% Dokumenttitel mit den Kopfangaben, darunter Empfänger und Absender
%% nebeneinander, unten der Fuß in drei Spalten. Die einzige Farbe ist die des
%% Logos oben links.
\documentclass[10pt,a4paper]{article}

%% pdfx erzeugt das PDF/A-1b, siehe balken.tex. Früh laden, es zieht hyperref
%% und xcolor mit.
\usepackage[a-1b]{pdfx}

\usepackage{fontspec}
\usepackage{geometry}
\usepackage{tabularx}
\usepackage{array}
\usepackage{etoolbox}  %% \ifblank für die Angaben, die leer sein dürfen
\usepackage{tikz}      %% das Logo

%% Fira Sans, aus TeX Live. Ligatures=TeXOff aus demselben Grund wie in
%% balken.tex: Text aus den Eingabedaten soll so dastehen, wie er in der
%% YAML-Datei steht, auch "--" und "<<" (#81).
\setmainfont{FiraSans}[
  Ligatures      = TeXOff,
  Extension      = .otf,
  UprightFont    = *-Regular,
  BoldFont       = *-Bold,
  ItalicFont     = *-Italic,
  BoldItalicFont = *-BoldItalic,
]

\geometry{a4paper, left=20mm, right=20mm, top=18mm, bottom=16mm}

\setlength{\parindent}{0pt}
\pagestyle{empty}
\raggedright

%% Keine Silbentrennung, aus demselben Grund wie in balken.tex: Eigennamen,
%% Beträge und Anschriften.
\hyphenpenalty=10000
\exhyphenpenalty=10000
\tolerance=2000

%% Gewöhnlicher Wortzwischenraum nach jedem Punkt (#143); den breiten
%% bekommen nur die eigenen Sätze des Programms in \hinweisabsatz.
\frenchspacing

%% Etwas Luft in den Zeilen des Gitters, sonst kleben die Ziffern an den Linien.
\renewcommand{\arraystretch}{1.35}

%% Die Breite der Betragsspalte. Positionen, Summen und Absetzung stehen in
%% drei Tabellen untereinander; nur mit derselben Breite fluchten ihre Linien.
%% Nachgemessen braucht "−19.040,00 €" rund 21 mm.
\newlength{\betragsbreite}\setlength{\betragsbreite}{24mm}
\newcolumntype{B}{>{\raggedleft\arraybackslash}p{\betragsbreite}}

%% Die Zeilen- und Tabellentrenner, die der Renderer aufruft (TemplateMacro).
%% Im Gitter trennt eine Linie die Positionen; die Summen stehen ohne Linie
%% untereinander, erst die Abschlusszeile setzt wieder eine.
\newcommand{\positionstrenner}{\\\hline}
\newcommand{\summentrenner}{\\}
\newcommand{\zeilenwechsel}{\\}
\newcommand{\zusammen}[1]{\mbox{#1}}

%% Eine Zeile oben rechts: Beschriftung und Wert als Tabellenzeile. Ist der Wert
%% leer, entfällt die ganze Zeile samt Beschriftung – sonst stünde auf einer
%% Stornorechnung "Zahlbar bis" ohne Datum.
\newcommand{\kopfangabe}[2]{\ifblank{#2}{}{#1 & #2\\}}

%% Die eigenen Sätze des Programms unter den Summen, wie in balken.tex: nur hier
%% der breite Abstand nach einem Satzende (#143), deshalb keine Eingabedaten.
\newcommand{\hinweisabsatz}[1]{\par\vspace{5mm}\noindent%
  {\nonfrenchspacing\sfcode`\:=1000 \sfcode`\;=1000 \sfcode`\,=1000 #1}}

\newcommand{\absetzungskopf}[1]{\par\vspace{6mm}\noindent #1\par\vspace{1.5mm}}

%% Die abgesetzten Abschläge im selben Gitter wie die Summen: Bezeichnung links,
%% Betrag in der Betragsspalte. Der Renderer liefert die Zeilen ohne Zeilenende
%% nach der letzten, das setzt diese Tabelle.
\newcommand{\absetzungstabelle}[1]{\par\noindent%
  \begin{tabularx}{\linewidth}{|>{\raggedright\arraybackslash}X|B|}\hline#1\\\hline\end{tabularx}\par}

%% Die Zusatzzeile unter einer Position oder einem Abschlag. \normalfont, weil
%% die Bezeichnung einer Position fett steht und die Zeile darin liegt.
\newcommand{\positionsdetail}[1]{\newline{\normalfont\small #1}}
\newcommand{\absetzungsdetail}[1]{\newline{\small #1}}

%% Eine Zeile der Positionstabelle, sechs Spalten wie das Gitter: laufende
%% Nummer, Bezeichnung (fett, die Beschreibung darunter), Menge, Einheit,
%% Einzelpreis, Gesamt.
\newcommand{\positionszeile}[6]{#1 & \textbf{#2} & #3 & #4 & #5 & #6}

\newcommand{\summenzeile}[2]{#1 & #2}

%% Die fette Zeile unter der Linie, mit der ein Betragsblock endet. Sie beendet
%% die Zeile davor selbst, wie in balken.tex (#147).
\newcommand{\abschlusszeile}[2]{\\\hline\textbf{#1} & \textbf{#2}}

%% Die Bankverbindung setzt raster nicht als eigenen Block: Sie steht im Fuß,
%% auf jedem Dokument, als Angabe des Briefbogens und nicht als
%% Zahlungsaufforderung – so, wie balken die IBAN in seiner Fußzeile führt.
%% Ob das Dokument zur Zahlung auffordert, sagt das Zahlungsziel oben rechts.
%% Der Renderer ruft \bankblock trotzdem auf, wenn Zahlung gefordert wird, und
%% es bleibt leer (#189, ADR-0011).
\newcommand{\bankblock}{}

%% Das Logo, fest in der Vorlage (#197): ein Beispiellogo der fiktiven
%% Muster IT-Consulting, eine Bildmarke "M" mit blauem Verlauf und eine
%% Wortmarke. Konfigurierbar wird es später.
%%
%% TikZ und keine Bilddatei: Der Treiber luatex.def von graphicx gibt jedem
%% über \includegraphics eingebundenen PDF eine Transparenzgruppe mit, und die
%% verbietet PDF/A-1 (ISO 19005-1, 6.4). TikZ-Code erzeugt keine. Außerdem
%% bleibt die Vorlage eine einzige Datei, und das Logo geht mit in die
%% Prüfsumme, aus der LatexCompiler die Trailer-ID bildet.
%%
%% Die Wortmarke hat ihre eigene Schrift: Ein Logo sieht gleich aus, egal in
%% welcher Schrift der Brief steht. Gezeichnet ist es in beliebigen Einheiten,
%% \resizebox bringt es auf \logohoehe.
\newfontfamily{\logoschrift}{texgyreheros}[
  Extension   = .otf,
  UprightFont = *-regular,
  BoldFont    = *-bold,
]
\definecolor{logodunkel}{HTML}{1B2A4A}
\definecolor{logohell}{HTML}{4FB3F0}
\definecolor{logomittel}{HTML}{2F7FD8}
\definecolor{logofalte}{HTML}{1F5FB0}
\newlength{\logohoehe}\setlength{\logohoehe}{12mm}
\newlength{\versalhoehe}
\newcommand{\logo}{\resizebox{!}{\logohoehe}{\begin{tikzpicture}
  %% Bildmarke: zwei Pfeiler und ein gefaltetes Band dazwischen
  \shade[top color=logohell, bottom color=logomittel, rounded corners=3pt]
    (0,0) -- (0,2.6) -- (0.8,2.6) -- (0.8,0) -- cycle;
  \shade[top color=logohell, bottom color=logomittel]
    (0.8,2.6) -- (1.6,1.5) -- (1.6,0.6) -- (0.8,1.7) -- cycle;
  \shade[top color=logohell, bottom color=logomittel]
    (1.6,1.5) -- (2.4,2.6) -- (2.4,1.7) -- (1.6,0.6) -- cycle;
  \shade[top color=logohell, bottom color=logomittel, rounded corners=3pt]
    (2.4,0) -- (2.4,2.6) -- (3.2,2.6) -- (3.2,0) -- cycle;
  \fill[logofalte] (0.8,2.6) -- (0.8,1.7) -- (1.15,2.12) -- cycle;
  \fill[logofalte] (2.4,2.6) -- (2.4,1.7) -- (2.05,2.12) -- cycle;
  %% Wortmarke
  \node[anchor=south west, inner sep=0, font=\logoschrift\fontsize{40}{40}\selectfont\bfseries]
    at (4.0,0.95) {\textcolor{logodunkel}{MUSTER}\textcolor{logohell}{IT}};
  \node[anchor=south west, inner sep=0, text=logodunkel,
    font=\logoschrift\fontsize{13}{13}\selectfont\addfontfeature{LetterSpace=30}]
    at (4.05,0.05) {CONSULTING};
\end{tikzpicture}}}

\begin{document}

%% ------------------------------------------------------- Titel und Kopfangaben
%% Links das Logo, rechts der Titel mit den Kopfangaben; die Beschriftungen
%% stehen rechtsbündig vor ihren Werten.
%%
%% Das Logo nimmt keinen Raum ein ([0pt][0pt]): So steht der Titel auf
%% derselben Grundlinie wie ohne Logo, und alles darunter bleibt, wo es war.
%% Seine Oberkante liegt auf der Versalhöhe des Titels. Dafür darf es nicht
%% höher werden als der Block der Kopfangaben, sonst ragt es in die Anschriften.
\settoheight{\versalhoehe}{\LARGE\bfseries M}%
\noindent\raisebox{\dimexpr\versalhoehe-\height\relax}[0pt][0pt]{\logo}%
\hfill
\begin{tabular}[t]{@{}r@{\hspace{3mm}}l@{}}
  \multicolumn{2}{@{}r@{}}{\LARGE\bfseries <<DOKUMENTART>>}\\[2mm]
  \kopfangabe{Rechnungsnr.:}{<<RECHNUNGSNUMMER>>}
  \kopfangabe{Kundennr.:}{<<KUNDENNUMMER>>}
  \kopfangabe{Unsere Auftragsnr.:}{<<AUFTRAGSNUMMER>>}
  \kopfangabe{Datum:}{<<RECHNUNGSDATUM>>}
  \kopfangabe{<<ZEITRAUM_BESCHRIFTUNG>>:}{<<LEISTUNGSZEITRAUM>>}
  \kopfangabe{Vereinnahmt am:}{<<VEREINNAHMT_AM>>}
  \kopfangabe{Ihre Referenz:}{<<KUNDENREFERENZ>>}
  \kopfangabe{Zahlbar bis:}{<<ZAHLUNGSZIEL>>}
\end{tabular}

\vspace{12mm}

%% ------------------------------------------------------ Empfänger und Absender
%% Links der Empfänger mit seiner USt-IdNr. (BT-48, §14a UStG), rechts der
%% Absender. Die Person unter der Firmierung nur ohne Geschäftsführung, wie in
%% balken; die Geschäftsführung steht hier im Fuß.
\begin{minipage}[t]{0.5\linewidth}\raggedright
  <<EMPFAENGER>>\ifblank{<<KUNDE_USTID>>}{}{\\[2mm]USt-IdNr.: <<KUNDE_USTID>>}
\end{minipage}%
\hfill
\begin{minipage}[t]{0.45\linewidth}\raggedleft
  <<FIRMA>>\ifblank{<<GESCHAEFTSFUEHRUNG>>}{\zeilenwechsel{}<<INHABER>>}{}\zeilenwechsel{}<<STRASSE>>\zeilenwechsel{}<<PLZ>> <<ORT>>\ifblank{<<TELEFON>>}{}{\zeilenwechsel{}Tel.: <<TELEFON>>}\zeilenwechsel{}<<EMAIL>>
\end{minipage}

\vspace{12mm}

%% --------------------------------------------------------------- Positionen
\noindent
\begin{tabularx}{\linewidth}{|>{\centering\arraybackslash}p{7mm}|>{\raggedright\arraybackslash}X|>{\raggedleft\arraybackslash}p{13mm}|>{\raggedright\arraybackslash}p{15mm}|>{\raggedleft\arraybackslash}p{22mm}|B|}
  \hline
  \textbf{Pos.} & \textbf{Bezeichnung} & \textbf{Menge} & \textbf{Einheit} & \textbf{Einzelpreis} & \textbf{Gesamt}\\
  \hline
<<POSITIONEN>>\\
  \hline
\end{tabularx}
%% ------------------------------------------------------------------ Summen
%% Direkt unter den Positionen, ohne Abstand: Mit derselben Betragsspalte
%% setzt sich das Gitter fort. \nointerlineskip nimmt den Zeilenabstand
%% zwischen den beiden Tabellen weg.
\par\nointerlineskip\noindent
\begin{tabularx}{\linewidth}{|>{\raggedright\arraybackslash}X|B|}
<<SUMMEN>>\\
  \hline
\end{tabularx}

%% ------------------------------------------------------ Die fachlichen Sätze
%% Pflichtsatz, Steuerhinweise, Absetzung und Skonto in derselben Reihenfolge
%% wie in balken. Jeder bringt seinen Abstand selbst mit.
<<ARTHINWEIS>>

<<STEUERHINWEISE>>

<<ABSETZUNG>><<SKONTO>>

<<BANKVERBINDUNG>>

\vspace{8mm}

%% Der Schlusssatz aus Stammdaten oder Rechnung. Keine Grußformel.
<<SCHLUSSSATZ>>

\vfill

%% ------------------------------------------------------------------ Fuß
%% Drei Spalten: Anschrift und Kontakt, Steuer und Register (§35a GmbHG),
%% Bank. Beide Steuerkennungen, wo es beide gibt – die Spalte hat Platz.
{\footnotesize
\noindent
\begin{minipage}[t]{0.32\linewidth}\raggedright
  <<FIRMA>>\zeilenwechsel{}<<STRASSE>>\zeilenwechsel{}<<PLZ>> <<ORT>>\ifblank{<<WEBSEITE>>}{}{\zeilenwechsel{}<<WEBSEITE>>}
\end{minipage}%
\hfill
\begin{minipage}[t]{0.34\linewidth}\raggedright
  USt-IdNr.: <<USTID>>\ifblank{<<STEUERNUMMER>>}{}{\\Steuernummer: <<STEUERNUMMER>>}\ifblank{<<REGISTERGERICHT>>}{}{\\Registergericht: <<REGISTERGERICHT>>\ifblank{<<REGISTERNUMMER>>}{}{, <<REGISTERNUMMER>>}}\ifblank{<<GESCHAEFTSFUEHRUNG>>}{}{\\Geschäftsführung: <<GESCHAEFTSFUEHRUNG>>}
\end{minipage}%
\hfill
\begin{minipage}[t]{0.28\linewidth}\raggedright
  <<INSTITUT>>\\IBAN: <<IBAN>>\\BIC: <<BIC>>
\end{minipage}\par}

\end{document}
